\chapter{Approaches I tested and did not use}
\label{ch:notused}

During the work I tested several ideas that did not become part of the final pipeline. I describe them
briefly here, with the reason why I did not use them.

\section{Making the simulations faster}
The OpenFOAM simulations are the slowest part of the pipeline: one setting of 600~s took 15--90~hours on one
CPU core of our lab computer, depending on the window speeds (faster air needs smaller time steps). The
\cotwo{} solver and the training run on the GPU: computing the \cotwo{} of one setting takes 1.25--4~hours, and
training the 5 networks about 1.5~hours. \tref{tab:speed} shows what I tried on the slowest setting (N12, all
three windows at 2~m/s).

\begin{table}[h]
  \centering
  \caption{Speed tests on setting N12.}
  \label{tab:speed}
  \begin{tabular}{>{\raggedright}p{5cm}rp{6.3cm}}
    \toprule
    Change & Speed-up & Result\\
    \midrule
    4 cores instead of 1 & 2.7$\times$ & identical results; used when only few settings run\\
    450~s instead of 600~s & 1.33$\times$ & \cotwo{} 1.2--1.5\,\% different (N10: 3.1\,\%)\\
    300~s instead of 600~s & 2$\times$ & about 4\,\% different\\
    Courant number 2 instead of 1 & 1.8$\times$ & flow about 20\,\% different; rejected\\
    looser solver tolerances & 1.1$\times$ & not used\\
    \bottomrule
  \end{tabular}
\end{table}

For many settings it is most efficient to run one setting per core; I use 4 cores only for single or slow
settings. A shorter simulated time of 450~s would save a quarter of the time with small differences, but I
kept 600~s for all settings so that the whole dataset is consistent. The larger Courant number changes the
flow too much, and looser tolerances hardly help.

\section{Running the simulations on the GPU}
To find out whether the simulations can be made much faster on the GPU, I wrote a lattice Boltzmann solver in
PyTorch (D3Q19, regularised BGK collision with a Smagorinsky model \citep{smagorinsky1963,kruger2017}). It uses
the same room, windows, door, columns and seating area. I tested setting N10 on a 4~cm lattice (7 million
nodes) for 600~s and compared it with OpenFOAM:
\begin{itemize}
  \item it took 7.3~hours on the RTX A2000, which is not faster than OpenFOAM on 4 cores;
  \item the mean flow differs from OpenFOAM by 68\,\% and the \cotwo{} at 600~s by 71\,\% (paper's metric
        11.1\,\%, while my \cotwo{} solver on the OpenFOAM flow gives 4.1\,\%).
\end{itemize}
The jet from window W1 slows down to 0.27~m/s within 2~m in OpenFOAM (RANS), but still has 0.47~m/s there in
the lattice Boltzmann solution (LES), which points to the different turbulence modelling. Without
measurements neither result can be called correct, and the dataset should stay consistent, so I did not use
this solver.

I also checked other software that could run the same turbulence model (transient k-$\omega$ SST with a
transported scalar) on the GPU \todo{add references for each}:
\begin{itemize}
  \item \textbf{OpenFOAM with GPU linear solvers} (AmgX, PETSc4FOAM, Ginkgo): only the pressure solver runs on
        the GPU, so the gain is limited.
  \item \textbf{OpenFOAM v2606}: the first official release with GPU offloading; still experimental, and it is
        not documented which solvers and turbulence models run on the GPU.
  \item \textbf{RapidCFD}: an OpenFOAM fork that runs fully on the GPU, but it is based on OpenFOAM 2.3.1 and
        CUDA 8 and could not be built with the current tools on our computer.
  \item \textbf{nekRS}: open source and GPU-based, but with a different discretisation and a large effort to
        set up.
  \item \textbf{Lattice Boltzmann codes} (e.g.\ FluidX3D): fast, but LES instead of RANS (see above).
  \item \textbf{ANSYS Fluent}: its GPU solver supports k-$\omega$ SST, transient runs and species transport, so
        it could reproduce our setup. It needs a commercial or academic licence; the free student version
        seems to be limited to about 512,000 cells, less than our mesh.
        \todo{confirm the limit with official documentation}
\end{itemize}
\todo{run time of the same short test on the DGX Spark and the M5 Max}

\section{Training without simulation data}
In this approach the network never sees simulation data \citep{raissi2019,li2021pino}. At random points in the
room it is told how much its prediction breaks the equations for airflow and \cotwo{}, and it is trained until
the equations are almost fulfilled. Some rules are built into the network so that they always hold exactly:
no air through the walls, the right amount of air through each window and door, the start from rest, and
\cotwo{} proportional to the number of people. The inputs are the 8 window speeds, the number of people and the
time (up to 2~minutes), and the output is the air velocity and the \cotwo{} everywhere in the room (a GNOT
network with about 322,000 parameters).

I compared every version with an OpenFOAM solution of the same equations, and also with a closed room, where
an independent finite-difference solution of the \cotwo{} is available. Over 25 versions I fixed several
problems (\tref{tab:pinnfix}).

\begin{table}[h]
  \centering
  \caption{Problems found in the physics-only approach.}
  \label{tab:pinnfix}
  \begin{tabular}{>{\raggedright}p{6.5cm}p{7.2cm}}
    \toprule
    Problem & What I did\\
    \midrule
    \cotwo{} stayed at zero & rescaled all quantities so that every term has a similar size\\
    air leaked through the walls & built the wall and window rules into the network\\
    the sharp edges of the incoming jets could not be learned & smoothed the built-in base flow inside the room\\
    the network found a too smooth flow that still fits the equations & gave the wall condition more weight\\
    \bottomrule
  \end{tabular}
\end{table}

In the closed room the \cotwo{} became reasonably accurate: about 13\,\% difference to the finite-difference
solution on a horizontal plane in the best version. With open windows the airflow did not: the velocity error
against OpenFOAM fell from 74\,\% to about 50\,\% and then stopped improving, and 30,000 or 50,000 training
steps made no difference. The total amount of \cotwo{} is about right, but it is in the wrong places, because
the predicted airflow is too slow. One training run took about 17~hours on the GPU. Such wrong but
equation-fitting solutions are a known problem of physics-only training \todo{reference}.

While testing I also found a bug in how the training points were chosen: the window speeds were tied to the
position of the points, so some windows were never trained as open, and some walls never saw air flowing along
them. It was present from about version 4 to version 15 and fixed in version 16, so conclusions about which
change helped before version 16 are less reliable. For our room this approach was not accurate enough, and
the rest of the thesis uses the CFD data.
